\section{Newman Theorem --- Func Disjointness Lower Bound --- Dual Hard Sample}
%%% generated by the blueprint pipeline from TCSlib.CommunicationComplexity.NewmanTheorem.FuncDisjointnessLowerBound.DualHardSample
\subsection*{Newman Theorem --- Func Disjointness Lower Bound --- Dual Hard Sample: Introduction}

This module provides the sample-level half of the Alice/Bob symmetry used in the randomized
lower bound for set disjointness (the ``WLOG'' step of [RY20, Thm.\ 6.13]): the involution
\texttt{CommunicationComplexity.Functions.Disjointness.RandomizedLowerBound.dualHardSample}
on hard samples, which swaps the two players' roles and reverses the coordinate order,
together with transport lemmas describing how bits, sets, special coordinates and windowed
conditioning vectors of a hard sample behave under it.  The declaration documented here is
the shared engine behind the windowed-vector duality lemmas.

\subsection*{Newman Theorem --- Func Disjointness Lower Bound --- Dual Hard Sample: Results}

\begin{theorem}[Reversal of a Windowed Bit Vector]
\lean{CommunicationComplexity.Functions.Disjointness.RandomizedLowerBound.xWindow_dualHardSample_eq_reverse}
\label{CommunicationComplexity.Functions.Disjointness.RandomizedLowerBound.xWindow_dualHardSample_eq_reverse}
\leanok
\uses{CommunicationComplexity.Functions.Disjointness.RandomizedLowerBound.HardSample, CommunicationComplexity.Functions.Disjointness.RandomizedLowerBound.dualHardSample, CommunicationComplexity.Functions.Disjointness.RandomizedLowerBound.reverseBoolVector, CommunicationComplexity.Functions.Disjointness.RandomizedLowerBound.xBit, CommunicationComplexity.Functions.Disjointness.RandomizedLowerBound.xBit_dualHardSample, CommunicationComplexity.Functions.Disjointness.RandomizedLowerBound.yBit}
\difficulty{4}
Let $n$ be a positive integer, let $\omega$ be a hard sample on $n$ coordinates with
special coordinate $T$, and let $P,Q$ be decidable relations on pairs of coordinates
such that, for every $i$, $P(i,\operatorname{rev}(T))$ holds exactly when
$Q(T,\operatorname{rev}(i))$ holds, where $\operatorname{rev}$ reverses coordinate
order. The vector that retains Alice’s bit at coordinate $i$ in the dual sample exactly
when $P(i,\operatorname{rev}(T))$ holds, and is false otherwise, equals the coordinate
reversal of the vector that retains Bob’s bit at coordinate $i$ in $\omega$ exactly
when $Q(T,i)$ holds, and is false otherwise.
\end{theorem}
